% !PRINTED 438
\weberpage{438}

%<<CONT
also \gesperrt{rational} durch $(k)$ und $(lk)$ ausdrückbar ist, und aus § 72
ergibt sich, daß $M$ eine \gesperrt{ganze algebraische Zahl} ist.

Diesen Satz wenden wir an auf je zwei aufeinanderfolgende
Glieder der Reihe (7) und erhalten, wenn wir % [sic] Der Druck nennt die Reihe (7); gemeint ist die auf der vorigen Seite als (10) bezeichnete Reihe.
\[
\omega_1 = \frac{\omega}{p}, \quad \omega_2 = \frac{\omega_1}{p}, \quad \dotsc, \quad \omega_\varepsilon = \frac{\omega_{\varepsilon-1}}{p}
\]
setzen, und wenn $\varphi$ eine durch die Klasse $l$ vollständig bestimmte
rationale Funktion bedeutet:
\[
\begin{aligned}
M &= \left(\frac{\eta(\omega_1)}{\eta(\omega)}\right)^{24} &&= \varphi(k,\, lk), \\
(2)\quad M_1 &= \left(\frac{\eta(\omega_2)}{\eta(\omega_1)}\right)^{24} &&= \varphi(lk,\, l^2k), \\
&\dotsc\dotsc\dotsc\dotsc\dotsc\dotsc\dotsc\dotsc\dotsc\dotsc\dotsc\dotsc \\
M_{\varepsilon-1} &= \left(\frac{\eta(\omega_\varepsilon)}{\eta(\omega_{\varepsilon-1})}\right)^{24} &&= \varphi(l^{\varepsilon-1}k,\, k).
\end{aligned}
\]
Durch Multiplikation aller dieser Gleichungen folgt:
\weq{3}{\left(\frac{\eta(\omega_\varepsilon)}{\eta(\omega)}\right)^{24} = \varphi(k,\, lk)\,\varphi(lk,\, l^2k)\,\dotsc\,\varphi(l^{\varepsilon-1}k,\, k).}
Nach § 120 kann aber die rechte Seite dieser Gleichung als
rationale Funktion $\Phi(k, lk)$ von $(k)$ und $(lk)$ (mit rationalen Zahlkoeffizienten)
dargestellt werden.
Andererseits ist, da $l$ zum Exponenten $\varepsilon$ gehört, $l^\varepsilon$ mit der
Hauptform, also $l^\varepsilon k$ mit $k$ und $\omega_\varepsilon$ mit $\omega$ äquivalent, also besteht
eine Gleichung:
\weq{4}{\omega_\varepsilon = \frac{\gamma + \delta\,\omega}{\alpha + \beta\,\omega}, \quad \alpha\delta - \beta\gamma = 1,}
und es ist $\eta(\omega_\varepsilon) : \eta(\omega)$ nach § 38 zu bestimmen. Genügt, wie
wir angenommen haben, $\omega$ der Gleichung
\weq{5}{a\,\omega^2 + b\,\omega + c\,p^\varepsilon = 0,}
so erhält man aus (4), da $\omega_\varepsilon = \omega/p^\varepsilon$ ist:
\[
\beta\,\omega^2 + (\alpha - p^\varepsilon\delta)\,\omega + \gamma\,p^\varepsilon = 0,
\]
was durch Vergleichung mit (5) gibt
\[
\beta = a\,x, \quad \alpha - p^\varepsilon\delta = b\,x, \quad \gamma = -c\,x,
\]
und wenn man $\alpha + p^\varepsilon\delta = y$ setzt:
\weq{6}{\alpha = \frac{y + b\,x}{2}, \quad \beta = a\,x,}
\[
p^\varepsilon\gamma = -c\,x, \quad p^\varepsilon\delta = \frac{y - b\,x}{2},
\]
%CONT>>
